%----------
%	LLM, AGENTES Y AGENT SKILLS PARA LA DETECCIÓN DE SESGO
%----------

En esa vía distinta, el conocimiento experto ya no se codifica en un diccionario,
sino que se comunica al modelo en lenguaje natural. Esta sección revisa esas tecnologías recientes (los LLM, las arquitecturas de agentes y las \textit{Agent Skills}), que
enmarca la propuesta de este trabajo.

\subsection{De los sistemas de reglas a los modelos de lenguaje}

Frente a los clasificadores supervisados clásicos, un LLM puede abordar la
detección de sesgo sin entrenamiento específico, guiado únicamente por
instrucciones. Este cambio es especialmente útil en lenguas con género gramatical
marcado, como el español, donde el fenómeno es más sutil que en inglés. Donde el
inglés emplea un término neutro como \textit{researchers}, el español obliga a
marcar el género, y el masculino genérico (``los investigadores'' para un equipo
liderado por mujeres) puede borrar la autoría femenina de forma casi imperceptible,
como en el caso de Bindi. El trabajo
de Derner et~al.~\cite{derner2024gender} es una de las primeras propuestas de usar
LLM para identificar y clasificar el sesgo de representación de género en corpus
en español, con una disparidad masculino:femenino que oscila entre 4:1 y 6:1 según
el corpus. En paralelo, han surgido marcos de evaluación pensados para el español y
su contexto cultural, como SESGO~\cite{robles2025sesgo}, que mide la tendencia de
los modelos a reproducir estereotipos en escenarios ambiguos, o el
\textit{toolkit} E.D.I.A.~\cite{edia2024} de la Fundación Vía Libre, que permite
auditar sesgos y estereotipos sin necesidad de conocimientos técnicos. Estos
trabajos confirman que el problema no desaparece al pasar a los LLM: los modelos
heredan el sesgo de sus datos, y detectarlo requiere método, no solo capacidad.

\subsection{Arquitecturas de agentes basadas en LLM}

Una \textit{agente} basada en LLM va más allá de una única llamada al modelo: combina
razonamiento, uso de herramientas y observación de los resultados en un bucle,
de modo que el modelo decide qué hacer a continuación en función de lo que
descubre. El patrón de referencia es ReAct~\cite{yao2023react}, que entrelaza
razonamiento y acción para reducir alucinaciones respecto a responder de un solo
tirón. Sobre esa base se han construido sistemas multiagente en los que varios
agentes especializadas se reparten una tarea compleja. En el ámbito periodístico,
AI-Press~\cite{liu2025aipress} asigna agentes a roles editoriales (recuperación,
redacción, revisión) apoyados en recuperación aumentada. Ahora bien,
automatizar parte del trabajo de las redacciones reactiva una exigencia clásica
cuando se delega en una máquina: la rendición de cuentas y la trazabilidad, es
decir, que ningún sistema opere como una caja negra y que toda afirmación pueda
rastrearse hasta su fuente~\cite{diakopoulos2019}. Esta exigencia es la que sostiene el
paradigma \textit{human-in-the-loop}: la agente no decide en lugar de la profesional,
sino que la asiste, y solo una salida trazable permite que la persona experta
audite, acepte o corrija cada resultado. Ambos principios, trazabilidad y
supervisión humana, guían el diseño del sistema propuesto en este trabajo.

\subsection{Agent Skills: conocimiento experto portable}

Un desarrollo reciente y directamente relevante para esta propuesta es el de las
\textit{Agent Skills}, formalizadas por Anthropic en 2025 como un formato abierto
para dotar a una agente de conocimiento especializado~\cite{anthropic2025skills}.
Una \textit{skill} es una carpeta con un fichero de metodología en lenguaje natural
y sus recursos asociados. Su mecanismo característico es la \textit{divulgación
progresiva} (\textit{progressive disclosure} en inglés): la agente solo ve el nombre y la
descripción de cada skill, y carga el cuerpo completo únicamente cuando la tarea lo
requiere, lo que evita saturar la ventana de contexto. Este principio no nace en el
ámbito de las agentes: procede de la disciplina de interacción persona-máquina (\textit{Human-Computer Interaction}, HCI por sus siglas en inglés), donde se
formuló como mostrar solo lo esencial y revelar el detalle bajo demanda, y se
popularizó en la literatura de usabilidad~\cite{nielsen1994}. El marco de
\textit{Agent Skills} le da una formulación operativa para modelos de lenguaje. La premisa que motiva
este enfoque es que buena parte del valor de un sistema reside en el conocimiento
experto que se le comunica, de modo que una metodología bien redactada y cargada
bajo demanda puede bastar para muchas tareas, sin recurrir a infraestructura más
pesada como un sistema de recuperación vectorial. Esta idea sitúa el conocimiento
experto en el centro del diseño y es la que articula la arquitectura del
Capítulo~\ref{implementation}, donde se contrasta de forma empírica frente a esas
alternativas.

\subsection{La intersección: agentes para detectar sesgo}

La combinación de agentes y detección de sesgo es la línea de investigación más reciente y donde
se sitúa este trabajo. Bias-Aware Agent~\cite{singh2025biasaware} incorpora la
detección de sesgo como una \textit{herramienta} de la agente en lugar de como un
paso posterior, y hace explícito qué fuentes se han usado y cuánto sesgo arrastra
cada una. Huang y Fan~\cite{huang2025fairness} proponen un sistema multiagente que
separa el hecho de la opinión antes de puntuar la intensidad del sesgo, y Huang y
Somasundaram~\cite{huang2024queer} muestran que una arquitectura de agentes
especializadas mejora de forma notable a un único modelo generalista en una tarea
de equidad lingüística. En el terreno específico del sesgo político y de género,
\textit{benchmarks} como EuroParlVote~\cite{yang2025europarlvote} evalúan estas
disparidades sobre datos parlamentarios europeos. Estos trabajos muestran que las agentes son una vía prometedora, pero ninguno lleva
la metodología experta del análisis de medios (los libros de códigos con
perspectiva de género) al interior de la agente ni la aplica al lenguaje sexista de la
prensa española. En ese cruce se sitúa este Trabajo de Fin de Máster, que encapsula
dicha metodología como \textit{agent skills} y evalúa de forma sistemática el valor real
de ese empaquetado.
